\فصل{مروری بر ادبیات و پیشینه پژوهش}

در این فصل، ادبیات مرتبط با مسئله زمان‌بندی کارگاهی به‌صورت نظام‌مند مرور و تحلیل می‌شود.
ابتدا رویکردهای کلاسیک حل، شامل روش‌های دقیق، قواعد ابتکاری
و الگوریتم‌های فراابتکاری، بررسی می‌شوند.
سپس پژوهش‌های فارسی مرتبط معرفی می‌شوند.
در ادامه، کارهای مرتبط با بهینه‌سازی چندهدفه در زمان‌بندی
و نیز زمان‌بندی در بستر رایانش لبه‌ای-مه-ابر مرور می‌شوند.
در پایان، جایگاه پژوهش حاضر نسبت به ادبیات موجود تبیین می‌شود.

\قسمت{رویکردهای حل مسئله زمان‌بندی کارگاهی}

مسئله زمان‌بندی کارگاهی از دهه ۱۹۶۰ مورد توجه پژوهشگران بوده است.
فیشر و تامپسون~\cite{fisher1963probabilistic} نخستین نمونه‌های معیار این مسئله را ارائه دادند.
آدامز و همکاران~\cite{adams1988shifting} روش گلوگاه متحرک\پاورقی{Shifting Bottleneck}
را برای حل دقیق نمونه‌های کوچک پیشنهاد کردند.
رویکردهای حل این مسئله را می‌توان به دسته‌های زیر تقسیم کرد:

\زیرقسمت{روش‌های دقیق}

روش‌های دقیق شامل شاخه و کران\پاورقی{Branch and Bound}
و برنامه‌ریزی خطی عدد صحیح\پاورقی{Integer Linear Programming} هستند.
اپلگیت و کوک~\cite{applegate1991computational}
مطالعه محاسباتی جامعی روی نمونه‌های معیار انجام دادند.
این روش‌ها عمدتاً برای نمونه‌های کوچک عملی‌اند.

\زیرقسمت{قواعد ابتکاری}

قواعد اولویت‌دهی\پاورقی{Dispatching Rules} از ساده‌ترین روش‌های حل‌اند.
دو قاعده پرکاربرد عبارتند از~\cite{pinedo2016scheduling}:

\شروع{فقرات}
\فقره \مهم{کوتاه‌ترین زمان پردازش} (\lr{SPT}\پاورقی{Shortest Processing Time}):
عملیاتی با کمترین زمان پردازش در اولویت قرار می‌گیرد.
\فقره \مهم{بیشترین کار باقی‌مانده} (\lr{MWR}\پاورقی{Most Work Remaining}):
عملیاتی که کار متناظر آن بیشترین زمان پردازش باقی‌مانده را دارد، در اولویت قرار می‌گیرد.
\پایان{فقرات}

این قواعد بسیار سریع‌اند؛ بااین‌حال، کیفیت راه‌حل‌های حاصل معمولاً در قیاس با روش‌های فراابتکاری پایین‌تر است.

\زیرقسمت{روش‌های فراابتکاری}

روش‌های فراابتکاری\پاورقی{Metaheuristic} از رایج‌ترین رویکردهای حل مسئله زمان‌بندی کارگاهی هستند.
مهم‌ترین این روش‌ها عبارتند از:

\شروع{فقرات}
\فقره \مهم{الگوریتم ژنتیک}:
هالند~\cite{holland1975adaptation} مبانی الگوریتم‌های ژنتیک را ارائه داد.
بیرویرث~\cite{bierwirth1995generalized} نمایش جایگشتی تعمیم‌یافته را برای مسئله زمان‌بندی کارگاهی پیشنهاد کرد.
چنگ و همکاران~\cite{cheng1996tutorial} مروری بر انواع نمایش‌ها در الگوریتم ژنتیک برای این مسئله ارائه دادند.
\فقره \مهم{تبرید شبیه‌سازی‌شده}:
این روش از فرآیند تبرید فلزات الهام گرفته شده و توانایی فرار از بهینه‌های محلی را دارد.
\فقره \مهم{بهینه‌سازی ازدحام ذرات}:
الگوریتمی مبتنی بر رفتار جمعی ذرات در فضای جست‌وجو.
\فقره \مهم{الگوریتم‌های میمتیک}\پاورقی{Memetic Algorithms}:
ترکیب الگوریتم تکاملی با جست‌وجوی محلی~\cite{moscato1989evolution}.
\پایان{فقرات}


\قسمت{ادبیات فارسی مسئله زمان‌بندی کارگاهی}

در ادبیات فارسی، پژوهش‌های متعددی درباره مسئله زمان‌بندی کارگاهی و مسائل مرتبط انجام شده‌اند.
در جدول~\رجوع{جدول:مقالات-فارسی} خلاصه‌ای از مهم‌ترین مقالات فارسی مرتبط ارائه شده است.

\شروع{لوح}[ht]
\تنظیم‌ازوسط
\شرح{خلاصه مقالات فارسی مرتبط با زمان‌بندی کارگاهی}

\شروع{جدول}{|c|r|r|}
\خط‌پر
\سیاه شناسه & \سیاه موضوع اصلی & \سیاه روش حل \\
\خط‌پر \خط‌پر
\lr{P1} & کارگاه باز منعطف & الگوریتم ژنتیک \\
\lr{P2} & زمان‌بندی کارگاهی کلاسیک & بهینه‌سازی ازدحام ذرات \\
\lr{P3} & زمان آماده‌سازی وابسته به توالی & مدل‌سازی ریاضی \\
\lr{P4} & مدل چندهدفه & تبرید شبیه‌سازی‌شده \\
\lr{P5} & کارگاهی منعطف با تعمیرات & فراابتکاری \\
\lr{P6} & زمان‌بندی بلادرنگ & یادگیری ماشین \\
\lr{P9} & کارگاهی منعطف با منابع دوگانه & فراابتکاری ترکیبی \\
\خط‌پر
\پایان{جدول}

\برچسب{جدول:مقالات-فارسی}
\پایان{لوح}

در ادامه، برخی از این مقالات با جزئیات بیشتری بررسی می‌شوند:

\مهم{حل مسئله زمان‌بندی در محیط کارگاه باز با الگوریتم ژنتیک}~\cite{persian_p1}:
این مقاله به حل مسئله زمان‌بندی در محیط کارگاه باز با ماشین‌های انعطاف‌پذیر با استفاده از الگوریتم ژنتیک می‌پردازد.
نمایش کروموزوم مبتنی بر جایگشت بوده و عملگرهای تقاطع و جهش خاص مسئله طراحی شده‌اند.

\مهم{به‌کارگیری الگوریتم ترکیبی بهینه‌سازی ازدحام ذرات}~\cite{persian_p2}:
در این پژوهش، الگوریتم ترکیبی بهینه‌سازی ازدحام ذرات برای حل مسئله سنتی زمان‌بندی کارگاهی ارائه شده است.
تابع هدف، کمینه‌سازی زمان تکمیل کل است.

\مهم{مدل ریاضی چندهدفه با تبرید شبیه‌سازی‌شده}~\cite{persian_p4}:
رحیمی و همکاران، یک مدل ریاضی چندهدفه برای زمان‌بندی در سامانه تولیدی کارگاهی طراحی کرده
و آن را با روش فراابتکاری تبرید شبیه‌سازی‌شده حل کرده‌اند.
اهداف شامل کمینه‌سازی زمان تکمیل کل و دیرکرد هستند.

\مهم{زمان‌بندی کارگاهی با زمان‌های آماده‌سازی}~\cite{persian_p3}:
اشقی و همکاران (دانشگاه صنعتی شریف) مدل‌سازی و حل مسئله زمان‌بندی کارگاهی
با زمان‌های آماده‌سازی وابسته به توالی را مطالعه کرده‌اند.

\مهم{زمان‌بندی بلادرنگ با یادگیری ماشین}~\cite{persian_p6}:
این پژوهش به زمان‌بندی کارگاهی بلادرنگ و نگهداری و تعمیرات پیشگویانه پویا
با استفاده از یادگیری ماشین پرداخته و رویکردی نوین ارائه کرده است.


\قسمت{بهینه‌سازی چندهدفه در زمان‌بندی}

بهینه‌سازی چندهدفه مسئله زمان‌بندی کارگاهی توجه فزاینده‌ای را به خود جلب کرده است.
دب و همکاران~\cite{deb2002fast} الگوریتم \lr{NSGA-II} را ارائه دادند
که به یکی از پرکاربردترین الگوریتم‌های چندهدفه تبدیل شده است.
کاسم و همکاران~\cite{kacem2002approach} از بهینه‌سازی تکاملی چندهدفه
برای مسئله زمان‌بندی کارگاهی منعطف استفاده کردند.

در حوزه شاخص‌های ارزیابی، زیتسلر و همکاران~\cite{zitzler2003performance}
مروری جامع بر شاخص‌های عملکرد الگوریتم‌های چندهدفه ارائه دادند.
شاخص ابرحجم\پاورقی{Hypervolume} به‌عنوان تنها شاخص سازگار با پارتو شناخته می‌شود~\cite{while2012fast}.


\قسمت{زمان‌بندی در بستر رایانش لبه‌ای، مه و ابر}

با گسترش اینترنت اشیاء و رایانش ابری، زمان‌بندی وظایف در زیرساخت‌های ناهمگن
به یک حوزه پژوهشی مهم تبدیل شده است.
شی و همکاران~\cite{shi2016edge} چالش‌ها و فرصت‌های رایانش لبه‌ای را مطرح کردند.
بونومی و همکاران~\cite{bonomi2012fog} مفهوم رایانش مه را معرفی نمودند.

لرا و همکاران~\cite{lera2019yafs} شبیه‌ساز \lr{YAFS} را برای شبیه‌سازی سناریوهای اینترنت اشیاء
در محیط رایانش مه ارائه دادند.
این شبیه‌ساز امکان مدل‌سازی توپولوژی شبکه، پهنای باند پیوندها
و توان پردازشی گره‌ها را فراهم می‌کند.


\قسمت{انتخاب تطبیقی عملگر}

انتخاب تطبیقی عملگر\پاورقی{Adaptive Operator Selection}
به انتخاب پویای عملگرهای ژنتیکی بر اساس عملکرد گذشته آن‌ها اشاره دارد.
آئر و همکاران~\cite{auer2002finite} الگوریتم \lr{UCB1} را
برای مسئله راهزن چندبازویی\پاورقی{Multi-Armed Bandit} ارائه دادند.
فیالهو و همکاران~\cite{fialho2010analyzing} کاربرد سازوکارهای مبتنی بر راهزن چندبازویی
در انتخاب عملگر را تحلیل کرده‌اند.

فرمول \lr{UCB1} برای انتخاب عملگر $i$ عبارت است از:
\begin{equation}
\text{UCB1}(i) = \bar{r}_i + C \sqrt{\frac{\ln N}{n_i}}
\end{equation}
که $\bar{r}_i$ پاداش میانگین عملگر $i$، $N$ تعداد کل انتخاب‌ها،
$n_i$ تعداد دفعات انتخاب عملگر $i$، و $C$ ضریب اکتشاف است.


\قسمت{جمع‌بندی و جایگاه پژوهش حاضر}

با مرور ادبیات، مشخص می‌شود که:

\شروع{شمارش}
\فقره اکثر پژوهش‌های فارسی بر حالت تک‌هدفه یا دوهدفه تمرکز دارند
و رویکرد سه‌هدفه در بستر لبه-مه-ابر کمتر مورد مطالعه قرار گرفته است.
\فقره ترکیب انتخاب تطبیقی عملگر با \lr{NSGA-II} برای مسئله زمان‌بندی کارگاهی
کمتر بررسی شده است.
\فقره شبیه‌سازی واقع‌گرایانه زیرساخت (با مدل انرژی و قابلیت اطمینان)
در کنار بهینه‌سازی زمان‌بندی، رویکردی نوین است.
\پایان{شمارش}

\شروع{لوح}[ht]
\تنظیم‌ازوسط
\شرح{مقایسه تطبیقی کارهای پیشین با پژوهش حاضر}

\شروع{جدول}{|c|c|c|c|c|c|c|}
\خط‌پر
\سیاه مرجع & \سیاه بستر & \سیاه تعداد اهداف & \سیاه روش & \سیاه انرژی/اطمینان & \سیاه \lr{AOS} & \سیاه حذفی \\
\خط‌پر \خط‌پر
\cite{persian_p1} & کارگاهی & ۱ & \lr{GA} & خیر & خیر & خیر \\
\cite{persian_p2} & کارگاهی & ۱ & \lr{PSO} هیبرید & خیر & خیر & خیر \\
\cite{persian_p4} & کارگاهی & ۲ & \lr{SA} & خیر & خیر & خیر \\
\cite{kacem2002approach} & \lr{FJSP} & ۲ & \lr{MOEA} & خیر & خیر & خیر \\
\cite{lera2019yafs} & مه/لبه & -- & شبیه‌سازی & جزئی & خیر & خیر \\
\مهم{پژوهش حاضر} & لبه-مه-ابر & ۳ & \lr{NSGA-II}+ & بله & بله & بله \\
\خط‌پر
\پایان{جدول}

\برچسب{جدول:مقایسه-کارهای-پیشین}
\پایان{لوح}

بر اساس جدول~\رجوع{جدول:مقایسه-کارهای-پیشین}، تمایز اصلی پژوهش حاضر در
ترکیب هم‌زمان چهار مؤلفه زیر است:
مدل‌سازی سه‌هدفه (زمان، انرژی و قابلیت اطمینان)،
در نظر گرفتن زیرساخت ناهمگن لبه-مه-ابر،
به‌کارگیری انتخاب تطبیقی عملگر در چارچوب \lr{NSGA-II}،
و انجام مطالعه حذفی برای سنجش سهم هر مؤلفه.
بدین‌ترتیب، نوآوری پژوهش صرفاً در بهبود یک عملگر خلاصه نمی‌شود، بلکه در یکپارچه‌سازی
مدل مسئله، روش حل و پروتکل ارزیابیِ قابل بازتولید نمود می‌یابد.
